\documentclass[11pt,a4paper]{article}
\usepackage[margin=1in]{geometry}
\usepackage{amsmath,amssymb}
\usepackage{booktabs}
\usepackage{graphicx}
\usepackage{hyperref}
\usepackage{listings}
\usepackage{xcolor}

\hypersetup{colorlinks=true,linkcolor=blue,citecolor=blue,urlcolor=blue}

\lstset{
  basicstyle=\ttfamily\small,
  breaklines=true,
  frame=single,
  backgroundcolor=\color{gray!10}
}

\title{\textbf{PolynomiaZ: Polar Coordinate Compression via\\Mathematical Function Fitting on Virtual Disk Platters}}
\author{Rory Toma}
\date{May 2026}

\begin{document}
\maketitle

\begin{abstract}
We present PolynomiaZ (PLTZ), a lossless compression codec that maps input data onto a virtual disk platter geometry and detects mathematical patterns within each track. Tracks matching a known pattern type are stored as compact function descriptors rather than raw bytes. A cross-track optimization pass further compresses groups of tracks with identical or linearly-varying parameters into single meta-descriptors. PLTZ achieves 5--12$\times$ better compression than zlib on linearly structured numerical data (firmware images, satellite telemetry, SPICE ephemeris kernels) while maintaining less than 1\% overhead on incompressible data. The codec supports adaptive sector sizing, multi-platter splitting, configurable raw fallback compression, and word-width-aware analysis for 16/32/64-bit numerical arrays.
\end{abstract}

\section{Introduction}

Standard compression algorithms (zlib, bzip2, LZMA) operate by identifying byte-level repetition through sliding windows, Burrows-Wheeler transforms, or dictionary coding. These approaches excel at data with repeated substrings---natural language text, markup, log files---but fundamentally cannot recognize that a byte sequence like $[0, 1, 2, 3, \ldots, 255]$ represents a linear function. They see ``no repeated substrings'' and store it nearly verbatim.

PolynomiaZ takes a different approach: it treats the input as data written to a virtual disk platter and asks, for each track, ``what mathematical function produced these bytes?'' When a function is found, the track is stored as a compact descriptor (e.g., \texttt{LINEAR(start=0, step=1)} in 4 bytes instead of 256 raw bytes).

This approach is not general-purpose. It targets a specific domain: data with known mathematical structure arising from physical processes, calibration sequences, address tables, and numerical computation.

\section{Architecture}

\subsection{Disk Platter Model}

Input data of length $N$ bytes is mapped onto a virtual disk with geometry $(T, S)$ where $T$ is the number of concentric tracks and $S$ is the number of sectors per track, such that $T \times S \geq N$. Each byte occupies a position $(r, \theta)$ in polar coordinates:
\begin{align}
r &= \lfloor i / S \rfloor \quad \text{(track index)} \\
\theta &= \frac{2\pi (i \bmod S)}{S} \quad \text{(angular position)}
\end{align}
where $i$ is the byte's sequential position in the input.

\subsection{Adaptive Geometry Selection}

The codec evaluates multiple sector sizes $S \in \{8, 16, 32, 64, 128, 256, 512\}$ and selects the geometry that produces the smallest compressed output. Different data benefits from different geometries: small sectors expose patterns in structured headers, while large sectors capture long arithmetic sequences.

\subsection{Multi-Platter Splitting}

After analysis, tracks are separated into two platters:
\begin{itemize}
\item \textbf{Platter 0}: Tracks with detected patterns (structured data) and meta-groups
\item \textbf{Platter 1}: Tracks with no pattern (raw data, optionally compressed with deflate)
\end{itemize}

\subsection{Cross-Track Meta-Descriptors}

After per-track analysis, groups of four or more tracks sharing the same pattern type are examined for parameter-level patterns. When found, the entire group is encoded as a single meta-descriptor:

\begin{itemize}
\item \textbf{AllSame} --- $N$ tracks with identical parameters stored once. Cost: $O(N)$ for indices + $O(1)$ for parameters, vs.\ $O(N \times p)$ individually.
\item \textbf{ConstLinear} --- $N$ CONST tracks whose values form an arithmetic sequence. Stored as start + step (2 bytes) instead of $N$ individual values.
\item \textbf{LinearStartsLinear} --- $N$ LINEAR tracks with the same step but linearly-incrementing start values. Stored as base\_start + start\_step + step (6 bytes).
\end{itemize}

This optimization is particularly effective for firmware images with large padding regions (many identical CONST tracks) and data with systematic ramps across tracks.

\section{Pattern Detection}

Each track is tested against pattern detectors in priority order. When both RLE and REPEAT match, the smaller encoding is chosen.

\subsection{Pattern Types}

\begin{table}[h]
\centering
\begin{tabular}{@{}llll@{}}
\toprule
\textbf{Tag} & \textbf{Pattern} & \textbf{Detection Criterion} & \textbf{Storage Cost} \\
\midrule
0 & CONST & All bytes identical & 1 byte \\
1 & LINEAR & $x_i = a + bi$ for all $i$ & 4 bytes \\
2 & RLE & Run-length encoding saves space & $2n_{\text{runs}}$ bytes \\
3 & REPEAT & Repeating sub-pattern of period $p < S/3$ & $p$ bytes \\
4 & DELTA & $\leq 8$ unique deltas, all in $[-128, 127]$ & $S+1$ bytes \\
5 & PERIODIC & FFT with $k$ significant components & $4 + 18k$ bytes \\
6 & POLY & Polynomial degree $\leq 8$ exact fit & $1 + 8(d+1)$ bytes \\
9 & LINEAR\_WIDE & Linear sequence in 16/32-bit words & 19 bytes \\
10 & LINEAR\_F64 & Linear sequence in IEEE 754 doubles & 18 bytes \\
11 & DELTA\_F64 & f64 values with f32-representable deltas & $10 + 4(n-1)$ bytes \\
7 & RAW & No pattern (fallback) & $S$ bytes \\
8 & RAW\_COMPRESSED & RAW with deflate applied & $3 + n_c$ bytes \\
\bottomrule
\end{tabular}
\caption{Pattern types and their storage costs. $S$ = sectors per track, $d$ = polynomial degree, $k$ = FFT components, $n_c$ = compressed size.}
\end{table}

\subsection{Wide-Word Analysis}

A key innovation is reinterpreting byte tracks as arrays of wider integer or floating-point types. A track of 128 bytes that appears random at the byte level may contain 32 linearly-incrementing 32-bit addresses:
\[
\texttt{[0x80000000, 0x80000020, 0x80000040, \ldots]}
\]
The LINEAR\_WIDE detector unpacks bytes as little-endian 16-bit or 32-bit unsigned integers and checks for arithmetic sequences. LINEAR\_F64 and DELTA\_F64 do the same for IEEE 754 double-precision values, with byte-exact reconstruction verification to guarantee lossless compression.

\subsection{Lossless Guarantee}

All pattern detectors verify byte-exact reconstruction before accepting a match. For floating-point patterns (LINEAR\_F64, DELTA\_F64), the reconstructed bytes are compared against the original---not just the numerical values---ensuring no precision loss from floating-point arithmetic.

\section{Binary Format}

\begin{lstlisting}
Header (12 bytes):
  [4] Magic "PLTZ"
  [4] Original data length (uint32 LE)
  [2] Sectors per track (uint16 LE)
  [2] Number of platters (uint16 LE)

Per platter:
  [2] Number of entries (uint16 LE)
  Per entry (individual track):
    [2] Track index (uint16 LE)
    [1] Pattern tag (uint8)
    [...] Pattern-specific payload
  Per entry (meta-group):
    [2] Sentinel 0xFFFF
    [1] META_TAG (12)
    [1] Inner pattern type
    [2] Track count
    [2*N] Track indices
    [1] Sub-tag (0=AllSame, 1=ConstLinear, 2=LinearStartsLinear)
    [...] Meta parameters
\end{lstlisting}

The format has a fixed 12-byte header overhead. Each individual track costs 3 bytes for its index and tag, plus the pattern-specific payload. Meta-groups amortize the per-track overhead across all member tracks.

\section{Experimental Results}

\subsection{Structured Data (PLTZ Advantage)}

\begin{table}[h]
\centering
\begin{tabular}{@{}lrrrr@{}}
\toprule
\textbf{Data} & \textbf{Original} & \textbf{PLTZ} & \textbf{zlib} & \textbf{Improvement} \\
\midrule
Linear ramp (byte) 4KB & 4,096 B & 126 B & 315 B & 2.5$\times$ \\
Linear f64 epochs 8KB & 8,192 B & 350 B & 1,881 B & 5.4$\times$ \\
Segment addresses (u32) 4KB & 4,096 B & 182 B & 2,244 B & 12.3$\times$ \\
Firmware (padding + headers) 4KB & 4,096 B & 65 B & 307 B & 4.7$\times$ \\
Constant fill 64KB & 65,536 B & 278 B & 84 B & 0.3$\times$ \\
Linear ramps 16KB & 16,384 B & 153 B & 408 B & 2.7$\times$ \\
Mixed structured + random 4KB & 4,096 B & 2,150 B & 2,381 B & 1.1$\times$ \\
\bottomrule
\end{tabular}
\caption{PLTZ vs.\ zlib on mathematically structured data. Cross-track meta-descriptors enabled.}
\end{table}

\subsection{Unstructured Data (PLTZ Disadvantage)}

\begin{table}[h]
\centering
\begin{tabular}{@{}lrrrl@{}}
\toprule
\textbf{Data} & \textbf{Original} & \textbf{PLTZ} & \textbf{zlib} & \textbf{Reason} \\
\midrule
English text & 1,580 B & 1,656 B & 680 B & Statistical, not mathematical \\
Random data & 4,096 B & 4,134 B & 4,109 B & Incompressible; 1\% overhead \\
Constant fill 64KB & 65,536 B & 278 B & 84 B & zlib RLE more byte-efficient \\
\bottomrule
\end{tabular}
\caption{Cases where standard compressors outperform PLTZ.}
\end{table}

\subsection{Domain-Specific Benchmarks}

\subsubsection{U-Boot Firmware}

Per-section analysis of a synthetic 512KB U-Boot image shows PLTZ winning on structured sections:
\begin{itemize}
\item Vector table: 36B vs.\ zlib 88B (2.4$\times$ better)
\item PLL configuration: 36B vs.\ zlib 201B (5.6$\times$ better)
\item Relocation table: 102B vs.\ zlib 576B (5.6$\times$ better)
\end{itemize}

\subsubsection{SPICE Ephemeris Kernels (JPL)}

For JPL SPICE kernel structures:
\begin{itemize}
\item Linear epoch arrays: 350B vs.\ zlib 1,881B (5.4$\times$ better)
\item Segment address directories: 182B vs.\ zlib 2,244B (12.3$\times$ better)
\item Chebyshev polynomial coefficients: PLTZ falls through to RAW; LZMA wins
\end{itemize}

\subsubsection{Cross-Track Impact}

The cross-track meta-descriptor optimization provides significant additional compression on data with many identical tracks:
\begin{itemize}
\item 64KB constant fill: 526B $\rightarrow$ 278B (1.9$\times$ improvement over per-track encoding)
\item 16KB linear ramps: 462B $\rightarrow$ 153B (3$\times$ improvement)
\item 4KB firmware image: 90B $\rightarrow$ 65B (1.4$\times$ improvement)
\end{itemize}

\section{Target Applications}

PLTZ is designed for domains where data originates from mathematical processes:

\begin{enumerate}
\item \textbf{Embedded firmware}: Flash images with 0xFF padding, address tables, calibration ramps
\item \textbf{Satellite telemetry}: Housekeeping data, attitude quaternions, epoch arrays, CCSDS framing
\item \textbf{Scientific instruments}: ADC/DAC sweeps, spectral calibration, sensor linearization tables
\item \textbf{SPICE kernels}: Epoch directories, segment address tables, orientation polynomials
\item \textbf{Protocol framing}: Sync patterns, incrementing sequence numbers, structured headers
\item \textbf{Lookup tables}: Gamma curves, PWM tables, motor control profiles
\end{enumerate}

\section{Implementation}

The primary implementation is written in Rust with a bzip2-style command-line interface:

\begin{lstlisting}[language=bash]
pltz file.bin           # compress -> file.bin.pltz
pltz -d file.bin.pltz   # decompress -> file.bin
pltz -k file.bin        # keep original
pltz -z file.bin        # enable deflate fallback for raw tracks
pltz -c file.bin        # output to stdout (pipeable)
pltz -t file.bin.pltz   # test integrity
\end{lstlisting}

The binary is approximately 1MB stripped, with no runtime dependencies beyond libc. The implementation includes 15 unit tests covering all pattern types, cross-track optimization, and edge cases.

\section{Comparison with Existing Work}

\begin{table}[h]
\centering
\begin{tabular}{@{}lll@{}}
\toprule
\textbf{Algorithm} & \textbf{Strategy} & \textbf{Best Domain} \\
\midrule
zlib (DEFLATE) & LZ77 + Huffman & Repeated substrings \\
bzip2 & BWT + RLE + Huffman & Long-range repetition \\
LZMA & Large dictionary + range coding & General purpose \\
\textbf{PLTZ} & \textbf{Function fitting on polar geometry} & \textbf{Mathematical structure} \\
\bottomrule
\end{tabular}
\caption{Algorithmic comparison.}
\end{table}

PLTZ occupies a unique niche: it recognizes \emph{semantic} structure (``this is a linear function'') rather than \emph{syntactic} structure (``these bytes appeared before''). The two approaches are complementary---PLTZ with a deflate fallback combines the strengths of both.

\section{Limitations and Future Work}

\textbf{Current limitations:}
\begin{itemize}
\item Single geometry for the entire file (mixed-structure files need pre-segmentation)
\item Periodic and polynomial detection not yet in Rust (Python reference only)
\item Slower than zlib due to adaptive geometry search (7 sector sizes $\times$ pattern detection)
\item Cross-track requires $\geq$4 tracks with same pattern to activate
\end{itemize}

\textbf{Planned improvements:}
\begin{itemize}
\item Columnar pre-processing transform for record-oriented data
\item Per-section geometry (different sector sizes for different regions)
\item Streaming/chunked mode for large files
\item File format versioning and magic number registry
\end{itemize}

\section{Conclusion}

PolynomiaZ demonstrates that domain-specific compression based on mathematical function fitting can dramatically outperform general-purpose algorithms on structured numerical data. By modeling data as points on a virtual disk platter, fitting functions per track, and compressing parameter metadata across tracks, PLTZ achieves 5--12$\times$ better compression than zlib on linear sequences, address tables, and epoch arrays---exactly the data structures found in firmware, telemetry, and scientific computing. The approach is complementary to existing compressors and is most effective when used as part of a hybrid pipeline or on data that is inherently mathematical in nature.

\end{document}
